\documentclass[aspectratio=169]{beamer}
\usetheme{default}

\usepackage{graphicx}
\usepackage{amsmath}
\usepackage{booktabs}

% UGA Colors
\definecolor{UGARed}{HTML}{BA0C2F}

% Minimal clean styling - only color frame titles
\setbeamercolor{frametitle}{fg=UGARed}
\setbeamercolor{title}{fg=UGARed}
\setbeamerfont{frametitle}{size=\Large, series=\bfseries}

% Remove navigation symbols
\setbeamertemplate{navigation symbols}{}

% UGA Logo in header and footer
\logo{\includegraphics[height=0.8cm]{assets/download.png}}

% Clean footer with UGA branding
\setbeamertemplate{footline}{
  \hbox{%
    \begin{beamercolorbox}[wd=\paperwidth,ht=2.5ex,dp=1ex,left]{footline}%
      \hspace{1em}\insertframenumber{}/\inserttotalframenumber
      \hfill\includegraphics[height=0.6cm]{assets/download.png}\hspace{1em}
    \end{beamercolorbox}%
  }
}

% Title info
\title{Mileage-Based Transportation Funding (VMT)}
\subtitle{Key Questions and Evidence}
\author{Presentation}
\date{February 2026}

\begin{document}

\frame{\titlepage}

% ============================================================
% SLIDE 1: The Three Questions
% ============================================================
\begin{frame}{Three Key Questions}
\begin{enumerate}
\item \textbf{Should transportation funding be based on how much you drive?}
\vspace{0.5cm}
\item \textbf{Would mileage-based fees disproportionately affect rural or low-income drivers?}
\vspace{0.5cm}
\item \textbf{How should privacy concerns be addressed if mileage tracking is required?}
\end{enumerate}
\end{frame}

%==============================================================================
% QUESTION 1: Should transportation funding be based on how much you drive?
%==============================================================================

% ============================================================
% SLIDE 2: Question 1 - Why VMT?
% ============================================================
\begin{frame}{Question 1: Should funding be based on miles driven?}
\textbf{Current Situation:}
\begin{itemize}
\item Average US Driver: 13--14K miles/year (FHA)
\item Fuel tax revenue declining
\item Oregon piloting VMT for 5,000 vehicles
\end{itemize}

\vspace{0.5cm}
\textbf{Key Advantage of VMT:}
\begin{itemize}
\item \textbf{All vehicles contribute} -- including electric vehicles (EVs)
\item Fuel tax doesn't capture EVs
\item More sustainable funding model
\end{itemize}
\end{frame}

% ============================================================
% SLIDE 3: Question 1 - State Variations
% ============================================================
\begin{frame}{Question 1: Driving Varies Dramatically by State}
\begin{columns}
\begin{column}{0.5\textwidth}
\textbf{Highest Miles Driven:}
\begin{itemize}
\item Wyoming: 21,589 miles
\item Indiana: 20,560 miles
\item Mississippi: 19,517 miles
\end{itemize}
\end{column}
\begin{column}{0.5\textwidth}
\textbf{Lowest Miles Driven:}
\begin{itemize}
\item Washington, D.C.: 6,694 miles
\item New York: 9,548 miles
\item Rhode Island: 9,903 miles
\end{itemize}
\end{column}
\end{columns}

\vspace{0.5cm}
\textbf{Georgia:} 17.5K miles/year (average)
\begin{itemize}
\item 2$\times$ more than NY resident
\item 3$\times$ more than D.C. resident
\end{itemize}

\vspace{0.3cm}
\textit{VMT tax better reflects actual road usage across states}
\end{frame}

%==============================================================================
% QUESTION 2: Would mileage-based fees disproportionately affect rural or low-income drivers?
%==============================================================================

% ============================================================
% SLIDE 4: Question 2 - Urban vs Rural
% ============================================================
\begin{frame}{Question 2: Impact on Rural vs. Urban Drivers}
\textbf{Common Claim:} VMT hurts rural residents

\vspace{0.5cm}
\textbf{Research Findings Say Otherwise:}
\begin{itemize}
\item Rural residents drive \textbf{more miles} than urban residents
\item Rural residents drive \textbf{less fuel-efficient} vehicles
\item Rural residents \textbf{already pay more} fuel tax
\item \textbf{VMT reduces the relative difference}
\end{itemize}

\vspace{0.5cm}
\begin{block}{Research Evidence}
McMullen et al. (2010): Rural households \textbf{benefit more} from VMT switch\\
Schroeckenthaler \& Fitzroy (2019): Urban areas pay \textbf{slightly more} under VMT
\end{block}
\end{frame}

% ============================================================
% SLIDE 5: Question 2 - Low vs High Income
% ============================================================
\begin{frame}{Question 2: Impact on Low-Income Drivers}
\textbf{Initial Concern:} VMT might burden low-income drivers
\begin{itemize}
\item Middle/upper class work remotely more
\item Blue-collar workers commute to work sites
\end{itemize}

\vspace{0.5cm}
\textbf{Research Findings:}
\begin{itemize}
\item VMT tax \textbf{is regressive} (larger \% of low-income budgets)
\item \textbf{BUT:} Fuel tax is \textbf{MORE regressive}
\item Higher-income households own newer, more fuel-efficient vehicles
\item \textbf{VMT reduces regressivity} compared to fuel tax
\end{itemize}

\vspace{0.3cm}
\begin{block}{Evidence}
Weatherford (2011): VMT shifts tax burden from \textbf{low to high income} households
\end{block}
\end{frame}

% ============================================================
% SLIDE 6: Question 2 - Public Perception vs Reality
% ============================================================
\begin{frame}{Question 2: Public Perception vs. Research}
\textbf{Public Perception (National Academies, 38 surveys):}
\begin{itemize}
\item Only 33--45\% say VMT is fair way to raise revenue
\item Only 15--38\% say VMT is fairer than fuel tax
\item Concerns: Not fair for rural, long-distance, work drivers
\end{itemize}

\vspace{0.5cm}
\begin{alertblock}{Reality Check}
``These opinions are \textbf{inconsistent with the empirical findings} regarding how VMT taxes actually vary across household categories.''
\end{alertblock}

\vspace{0.3cm}
\textbf{Support increases with understanding:}
\begin{itemize}
\item General public: 24\% support
\item Pilot program participants: 51\% support
\end{itemize}
\end{frame}

%==============================================================================
% QUESTION 3: How should privacy concerns be addressed if mileage tracking is required?
%==============================================================================

% ============================================================
% SLIDE 7: Question 3 - Privacy Concerns
% ============================================================
\begin{frame}{Question 3: Privacy Concerns Are Real}
\textbf{Key Findings (National Academies, 2016):}
\begin{itemize}
\item Only 4 of 38 surveys asked about privacy
\item \textbf{50\%+ of respondents} concerned about location tracking
\item Main objection: \textbf{Where and when} you drive
\item \textbf{Odometer-only data} did not generate undue concern
\end{itemize}

\vspace{0.5cm}
\textbf{Key Insight:}
\begin{itemize}
\item Location/time tracking = major concern
\item Simple odometer reading = acceptable
\end{itemize}
\end{frame}

% ============================================================
% SLIDE 8: Question 3 - Privacy Solutions
% ============================================================
\begin{frame}{Question 3: Addressing Privacy Concerns}
\textbf{Proposed Solutions:}
\begin{itemize}
\item \textbf{Use odometer-only collection} (no location/time data)
\item \textbf{Private firm collects data} (not government directly)
\item Minimize data collection to essential information only
\end{itemize}

\vspace{0.5cm}
\textbf{Context: We Already Accept Tracking:}
\begin{itemize}
\item E-ZPass: Tracks toll plaza passages
\item OnStar (GM): Vehicle location and driving data
\item Insurance apps: Allstate DriveWise, Progressive Snapshot
\item Cell phones: Location tracking by private and public entities
\end{itemize}

\vspace{0.3cm}
\textit{Privacy concerns must be addressed, but technology exists to do so}
\end{frame}

%==============================================================================
% CONCLUSION: Answering the Three Questions
%==============================================================================

% ============================================================
% SLIDE 9: Summary - Answering the Three Questions
% ============================================================
\begin{frame}{Answering the Three Questions}
\begin{block}{Q1: Should transportation funding be based on miles driven?}
\textbf{Yes} -- Includes all vehicles (EVs), sustainable funding, reflects actual usage
\end{block}

\begin{block}{Q2: Would this disproportionately affect rural or low-income drivers?}
\textbf{No} -- VMT \textbf{reduces} inequity compared to fuel tax (contrary to public perception)
\end{block}

\begin{block}{Q3: How should privacy concerns be addressed?}
\textbf{Odometer-only tracking + private collection} -- No location/time data needed
\end{block}

\vspace{0.3cm}
\textbf{Key Challenge:} Public perception doesn't match research findings\\
\textit{Education and transparency are critical}
\end{frame}

% ============================================================
% SLIDE 10: References
% ============================================================
\begin{frame}{References}
\small
\begin{enumerate}
\item David L. Sjoquist (2023). \textit{Funding Transportation in Georgia: Vehicle Miles Traveled Tax}. Center for State and Local Finance, Andrew Young School of Policy Studies, Georgia State University.

\item National Academies of Sciences, Engineering, and Medicine (2016). Analysis of public opinion surveys on mileage-based user fees. Cited in Sjoquist (2023).

\item Congressional Budget Office (2019). Reports on national vehicle miles traveled tax. Cited in Sjoquist (2023).

\item Government Accountability Office (2012). Reports on mileage-based user fees. Cited in Sjoquist (2023).

\item National Surface Transportation Infrastructure Financing Commission (2009). Reports on transportation funding mechanisms. Cited in Sjoquist (2023).

\item McMullen, B. S., Zhang, L., and Nakahara, K. (2010). Distributional impacts of mileage-based user fees. Cited in Sjoquist (2023).

\item Schroeckenthaler, J., and Fitzroy, S. (2019). Equity effects of mileage-based road usage charges. Cited in Sjoquist (2023).

\item Weatherford, B. A. (2011). Distributional impacts of a vehicle miles traveled tax. Cited in Sjoquist (2023).
\end{enumerate}
\end{frame}

\end{document}
